\subsection{五个分问的思路怎样连起来}

\begin{enumerate}
  \item 问题一先确定龙头在阿基米德螺线上的位置，相邻把手随后按距离不变的关系逐节递推；
  速度由短时间内的弧长变化得到。这一问最后既留下数值，也留下后四问共用的位置、速度计算器。
  \item 问题二把每块板凳写成矩形，用分离轴判断重叠。这里没有一上来就缩小检查范围，
  而是先用反证法说明第一次碰撞对象，再让程序在第一次触发处停止。
  \item 到了问题三，“能进入调头区域”被落实为全过程均不碰撞。螺距从粗区间缩到细步长，
  终点可行却中途碰撞的候选被舍去，所以可行性不是只看最后一帧。
  \item 问题四先由相切关系写出两段圆弧，再用“后续构件不能倒退”给大圆弧半径下界；
  确定路径后，位置递推按盘入、两段圆弧和盘出四种状态继续运行。
  \item 问题五不另建运动模型，只把问题四所得的各把手速度改写成上限约束；粒子群在这里只负责
  搜索龙头速度。这样，算法承担的任务很窄，模型链的职责也随之分明。
\end{enumerate}

